% =====================================================================
% 05-optimization.tex
% =====================================================================
\section{The Optimization Algorithm}\label{sec:optimization}

The second axis of our taxonomy concerns how the policy is updated once a
reward is available. The dominant methods form a spectrum from general
policy gradients to specialized preference learning; their trade-offs are
summarized in Table~\ref{tab:algos}.

\subsection{Policy Gradient with a Critic: PPO}

PPO~\cite{schulman2017ppo} optimizes a clipped surrogate that prevents
destructively large updates, using a learned value function to estimate
advantages $A(s_t,a_t)=R(s_t,a_t)-V(s_t)$. It was the engine of
InstructGPT and remains a strong baseline: a careful study found that
well-tuned PPO still outperforms DPO in complex settings that require
exploration, provided the reward model, KL penalty, and batch size are
well chosen~\cite{xu2024dpoppo}. PPO's principal cost is the critic
network, of size comparable to the policy, which roughly doubles memory
and complicates training.

\subsection{Critic-Free and Group-Relative Methods}

A second family removes the critic to reduce cost and instability.
REINFORCE with a simple baseline is the classical
instance~\cite{williams1992reinforce}, and RLOO revisits it for LLMs,
using a leave-one-out average over sampled responses as a low-variance
baseline that matches or beats PPO at lower
cost~\cite{ahmadian2024rloo}. GRPO~\cite{shao2024deepseekmath} samples a
\emph{group} of responses per prompt and normalizes their rewards within
the group to form advantages, dispensing with the value network
altogether. GRPO is especially effective with verifiable binary rewards
and underpins DeepSeek-R1~\cite{guo2025deepseekr1}; subsequent analysis
characterizes its effective loss and success-amplification
dynamics~\cite{shao2025grpodynamics}, and variants further relax the
reference-model and KL requirements for stability.

\subsection{Direct and Offline Preference Optimization}

The third family, anchored by DPO~\cite{rafailov2023dpo} and its
relatives IPO~\cite{azar2024ipo} and KTO~\cite{ethayarajh2024kto},
optimizes a closed-form preference loss directly, with no on-policy
sampling and no reward model. These methods are simple and stable but
fundamentally offline, and a growing body of evidence shows that online
methods converge to better final performance because they can explore
beyond the fixed preference
distribution~\cite{tajwar2024preference,xu2024dpoppo}. This motivates
hybrid schemes such as Nash learning from human
feedback~\cite{munos2023nash} and asynchronous off-policy
RLHF~\cite{noukhovitch2025async}, which seek the stability of offline
methods with the ceiling of online ones.

\subsection{Online Versus Offline: The Central Trade-Off}

The recurring theme is the online--offline trade-off. Offline methods
(DPO family) are cheap and stable but capped by their dataset; online
methods (PPO, GRPO) explore and reach higher performance but are
computationally expensive because every update requires fresh sampling
from the full model. Critic-free online methods (GRPO, RLOO) occupy an
attractive middle ground, which largely explains their current
popularity for reasoning-oriented training.

\begin{table}[t]
\centering
\caption{Comparison of the dominant optimization algorithms.}
\label{tab:algos}
\renewcommand{\arraystretch}{1.25}
\footnotesize
\begin{tabular}{@{}p{0.18\columnwidth}p{0.16\columnwidth}p{0.13\columnwidth}p{0.13\columnwidth}p{0.20\columnwidth}@{}}
\toprule
\textbf{Method} & \textbf{Reward model} & \textbf{Critic} & \textbf{Data} & \textbf{Best for} \\
\midrule
DPO  & Implicit & No  & Offline & General alignment \\
PPO  & Explicit & Yes & Online  & Complex / exploratory tasks \\
GRPO & Explicit (verifiable) & No & Online (groups) & Verifiable reasoning \\
RLOO & Explicit & No & Online & Low-variance, cheap \\
\bottomrule
\end{tabular}
\end{table}
